\section{Resultados obtenidos}

\subsection{Optimización de Procesos}

Con la implementación de Odoo, la agencia ha logrado estructurar un flujo operativo completamente integrado que conecta la gestión de prospectos, la venta de paquetes turísticos y la retroalimentación del cliente en un mismo ecosistema. A diferencia del modelo anterior, basado en la fragmentación de datos en Excel y WhatsApp, el nuevo flujo permite una trazabilidad continua:

\begin{itemize}
    \item \textbf{Gestión de Prospectos (CRM):} Cada consulta recibida por canales digitales se registra como una oportunidad comercial, permitiendo un seguimiento formal que antes dependía de la memoria o de chats dispersos.
    \item \textbf{Cotización Automatizada:} Al contar con un catálogo digital de servicios (Cusco, Arequipa, Puno), el sistema genera itinerarios y presupuestos profesionales en PDF de manera instantánea, eliminando la necesidad de digitar manualmente cada propuesta en Word.
    \item \textbf{Conversión a Orden de Venta:} Una vez aceptada la propuesta, la cotización se convierte en una Orden de Venta con un solo clic, asegurando que los datos del pasajero y las condiciones de pago (adelantos y saldos) se mantengan íntegros sin errores de transcripción.
    \item \textbf{Ciclo de Retroalimentación:} Tras la finalización del servicio, el sistema gestiona el envío de encuestas de satisfacción para capturar el "feedback" del turista, alimentando la reputación digital de la agencia.
\end{itemize}

\subsection{Impacto de la Productividad}

La digitalización ha generado una mejora sustancial en la eficiencia operativa de la agencia, reduciendo drásticamente los tiempos administrativos y los riesgos de error. Para validar estos resultados, se comparan los indicadores del diagnóstico inicial con el estado actual del sistema:

\begin{table}[htbp]
\centering
\renewcommand{\arraystretch}{1.25}
\setlength{\tabcolsep}{5pt}
\small
\begin{tabularx}{\textwidth}{|>{\raggedright\arraybackslash}p{3.0cm}|>{\raggedright\arraybackslash}X|>{\raggedright\arraybackslash}p{2.6cm}|>{\raggedright\arraybackslash}p{2.6cm}|}
\hline
\textbf{Indicador} & \textbf{Fórmula de Cálculo} & \textbf{Valor Inicial} & \textbf{Valor Final} \\ \hline
Tiempo de Cotización & $\frac{\sum Tiempo}{N^\circ Cotizaciones}$ & 45 - 60 minutos & $< 15$ minutos \\ \hline
Tasa de Error Operativo & $\left(\frac{N^\circ Errores}{N^\circ Reservas}\right)\times 100$ & 5\% & $< 1\%$ \\ \hline
Exactitud de Información & \% Expedientes 100\% digitales & "Silos" (Excel / WhatsApp) & 100\% en CRM \\ \hline
Satisfacción del Cliente & Mejora porcentual en NPS & Línea base inicial & +15\% (Meta) \\ \hline
\end{tabularx}
\normalsize
\end{table}

\subsection{Adopción del Sistema}

La adopción de Odoo se ha desarrollado de manera progresiva, centrada en la Sra. Silveria Flores y su equipo operativo. La transición representó un cambio cultural: de una gestión basada en la "flexibilidad manual" y la memoria, a un modelo de estandarización digital. Entre los hitos de adopción destacan:

\begin{itemize}
    \item Uso activo de la aplicación móvil para consultas rápidas de itinerarios y datos de pasajeros durante los viajes.
    \item Reducción drástica del uso de cuadernos y archivos Excel externos para el registro de clientes.
    \item Dominio del flujo de cotización-confirmación-venta dentro de la plataforma.
\end{itemize}

\subsection{Plan de Capacitaciones}

Se ejecutó un programa de transferencia de conocimiento para asegurar la sostenibilidad del sistema:

\begin{table}[htbp]
\centering
\renewcommand{\arraystretch}{1.25}
\begin{tabular}{|p{3.2cm}|p{7.8cm}|p{1.8cm}|}
\hline
\textbf{Área / Perfil} & \textbf{Tema Impartido} & \textbf{Horas} \\ \hline
Gerencia y Ventas & Introducción a Odoo y Navegación CRM & 2 \\ \hline
Operaciones & Configuración de Paquetes y Cotizador & 3 \\ \hline
Atención al Cliente & Gestión de Encuestas y Feedback & 2 \\ \hline
Tecnología & Uso de Aplicación Móvil y Seguridad & 2 \\ \hline
\end{tabular}
\end{table}
